\chapter{A five-root local certificate on an A-unbounded weak-C2 subchart}
\label{component:p9ec3d6bbd717}
\noindent\textbf{Component }\path{p9ec3d6bbd717}\par
{\small\noindent Source: \path{R17/definitions/weak_c2_complete_WA_frame.tex}\par}
{\small\noindent SHA256: \path{9ec3d6bbd7170be3f5d3cc22ae2d244ff93c9445f41550c5c9369804dab0bbcd}\par}
{\small\noindent Inventory use: WA3-WA8 and WA10\par}
\medskip
\section{A five-root local certificate on an A-unbounded weak-C2 subchart}
\label{p9ec3d6bbd717:sec:r9-weak-A-five-root}

\paragraph{Revision 10 scope update.}
The inherited boundary extension in \path{inherited_r16/Q222_Revision13/advances/weak_A_boundary/extension.tex} extends this admission to $z/B$
positive infinitesimal with all other source hypotheses retained.
The unit-ratio restriction below records the original revision-9 theorem.

\begin{theorem}
Retain the original physical C2 source, its exact stationary contact,
all strict source gaps, and $n\ge0$.
Normalize $d=1$. Suppose $B,z,y$ are positive infinitesimals,
$r,G/B,j,y/B,L/B$ are positive bounded units, and $H/y$ is
infinitesimal. Impose in addition
\[
 z/B\text{ and }1-r\text{ are positive bounded units}.       \tag{WA1}
\]
Then the entire original parent is a finite tensor sum of squares
for arbitrary independent testing half-angles $u,v$ that are
infinitesimal or zero, including either sign and complex states.
There is no comparison of a testing scale with $B,H,z,y$.
The total-scale parameter $\delta_{\rm tot}$ is positive infinitesimal.
This is a local testing certificate on the stated source subchart;
it does not assert that its intersection with the final residual
complement is nonempty, or supply an unproved remote-testing join.
\end{theorem}

No membership in the residual complement is needed in this proof.
When applied to the original residual problem, its existing complement
condition is still intersected with the source chart; no nonemptiness
of that intersection is inferred.

The original omitted event coefficients remain $\lambda=\beta=0$.
Here $B=H+L$, and the actual contact in the positive gauge is
\[
 (u_c,v_c;W_c,X_c,Y_c,Z_c)
 =\left(\sqrt z,\sqrt y;\frac B{\sqrt{zy}},
                  -r\sqrt z,\frac H{\sqrt y},1\right).
\]
The source ratio $\ell=B/z$ below is not an event coefficient.

\paragraph{The source pivots survive the small H opening.}
Write $h=H/B$, $\kappa=B/y$, $\ell=B/z$. Both $\kappa$ and
$\ell$ are positive bounded units, and $h$ is infinitesimal.
The original formulas imply that $n,m,C,D,\alpha$ are bounded,
that $m,C,D$ are positive bounded units, and that $Ah$ is a
positive bounded unit. Thus $A$ is unbounded and
$Ah^2,\sqrt A\,h,A^{-1/2}$ are infinitesimal. No bound on the
relative rates of $h$ and $B$ is used.

Here is an explicit check of the radial pivot. With $T_c=B+zr$,
the original formula is
\[
 \frac\alpha{B^2}
   =\frac{T_c}{\delta}\frac{M/B^2}{2(1-H)},\qquad
 M-yG=H(1+z)T_c+z\{yj+rL(1+y)\}>0.                  \tag{WA2}
\]
Consequently $M/B^2\ge(y/B)(G/B)$ is bounded below by a
positive real, and it is bounded above by the original formula
for $M$ and the unit ratios in (WA1). Since $T_c/\delta$ is a
positive bounded unit, $\alpha/B^2$ is a positive bounded unit
uniformly as $H/y\to0$. Likewise
\[
 m=\frac{z^2(1+y)(1-r)(1+zr)(1+zr+L)}{2G\delta}
\]
is a positive bounded unit under (WA1); the local residue argument
uses only $n\ge0$. Finally $A$ is the only unbounded retained coefficient, so
$\delta_{\rm tot}=1/(1+n+m+C+A+D+\alpha)$ is positive
infinitesimal. These arguments preserve all original coefficient
and source-gap hypotheses.

\paragraph{Five homogeneous contact roots.}
Put $U=u^2$, $V=v^2$, and define
\[
 \xi=X+ruZ,\qquad q=\sqrt A\,(Y-huW),\qquad
 R_1=uW-\kappa vZ,\quad R_2=vW-\ell uZ,\quad
 T=uvW-BZ.                                             \tag{WA3}
\]
Every root vanishes at the actual source contact. Inverses and
radicals in their coefficients are pure source. The large $A$
packet is retained in $q$; no positive packet is subtracted from
the parent and no physicality hypothesis is transferred to a
modified remainder.

Use the following source coefficients for the exact scalar graph:
\begin{align*}
 a&=(C+D+\tfrac12)(1-h)^2+Ah^2,&a_4&=Ah^2,\\
 b_2&=2mh+C(1-h)^2+(A+\alpha+\tfrac12)h^2,
 &b_4&=h^2(m+A+\alpha+\tfrac12),\\
 c&=\alpha,&c_2&=nr^2+(m+\tfrac12)(1-r)^2+\alpha,\\
 c_4&=nr^2,&d_2&=(n+\tfrac12)r^2+2Dr,\\
 d_4&=(n+D+\tfrac12)r^2.
\end{align*}
Thus the graph on $\xi=q=0$ is
\[
 \begin{split}
 F_g/4={}&[aU+a_4U^2+(m+b_2U+b_4U^2)V]W^2\\
 &+[c+c_2U+c_4U^2+(D+d_2U+d_4U^2)V]Z^2
       -2uv(e_0+e_1U)WZ,
 \end{split}
\]
where $e_0,e_1$ are exactly the original coefficients in (WS3).
Set
\[
 \begin{gathered}
 \beta_s=\frac{Bd_4+\kappa c_4}{\kappa\ell^2},\qquad
 \gamma_s=b_4+\frac{d_4}{\kappa\ell},\\
 j_0=\frac{Bd_4-a_4\kappa^2\ell+d_2\ell}
                   {2\kappa\ell^2},\qquad
 j_1=\frac{d_4}{2\kappa\ell},\qquad J_s(U)=j_0+j_1U,\\
 f_s=\frac12\left(b_2-\frac{c_4}{\ell^2}
                -\frac c{B^2}-\frac{Bd_4}{\kappa\ell^2}\right).
 \end{gathered}                                        \tag{WA4}
\]
The four original scalar contact equations (value and the three
logarithmic derivatives in $u,v,W/Z$) give the exact identity
\[
 F_g/4=(R_1,R_2,T)\,\mathsf S\,(R_1,R_2,T)^T,\qquad
 \mathsf S=
 \begin{pmatrix}
 a+a_4U&uvJ_s&-vJ_s\\
 uvJ_s&m+\beta_sU&u(f_s-j_1U)\\
 -vJ_s&u(f_s-j_1U)&c/B^2+\gamma_sU
 \end{pmatrix}.                                       \tag{WA5}
\]
This is a different scalar identity from the excluded polynomial
two-block splitter (JG7). In particular, the last diagonal retains
the full radial mass $\alpha/B^2$ and the Gram depends on the
independent testing variables. All source coefficients in (WA4)
are bounded; no uncancelled $1/B$ occurs.

\paragraph{Both full first normals and the complete hard block.}
For the $\xi$ normal define
\[
 \begin{gathered}
 p_0=\frac{D+r(n+1/2)}\kappa,\qquad
 p_1=\frac{r(n+D+1/2)}\kappa,\qquad
 q_1=\frac{nr+Bp_1}\ell,\\
 f_x=-mh+(D+\tfrac12)(1-h)-p_0-q_1,\\
 \mathsf L_x=\bigl(uv(p_0+p_1U),\ -m+q_1U,\
                                u(f_x-p_1U)\bigr).
 \end{gathered}                                        \tag{WA6}
\]
For the unnormalized $\eta=Y-huW$ normal set
\[
 \begin{gathered}
 p_y=-(C+D+\tfrac12)(1-h)+Ah,\qquad p_y'=Ah,\\
 k_v=-C(1-h)+Ah+m+(\alpha+\tfrac12)h,\qquad
 k_{uv}=h(A+m+\alpha+\tfrac12),\\
 k_z=-m(1-r)+r(D+\tfrac12)-\alpha-\tfrac12,\\
 q_y=\frac{-k_z-\kappa p_y'-Bk_{uv}}\ell,\qquad
 f_y=k_v-q_y,\\
 \mathsf L_y=\bigl(p_y+p_y'U,\ uvq_y,\
                                      v(f_y+k_{uv}U)\bigr).
 \end{gathered}                                        \tag{WA7}
\]
The two remaining original source stationarity equations, in $X,Y$,
give these exact representations of the complete first normals.
No first-normal term is discarded. Define also
\[
 \begin{split}
 h_{xx}&=n(1+U)(1+V)+m+(D+\tfrac12)UV+\tfrac12(1+V),\\
 h_{xy}&=-uv(m+D+\tfrac12),\\
 h_{yy}&=A(1+U)(1+V)+C(1+V)+D+\tfrac12
                   +mUV+(\alpha+\tfrac12)(1+U)V.
 \end{split}
\]
Then the \emph{entire original four-state parent} is exactly
\[
 F/4=\mathcal R^T\mathsf G(u,v)\mathcal R,\qquad
 \mathcal R=(\xi,q,R_1,R_2,T)^T,
\]
\[
 \mathsf G=
 \begin{pmatrix}
 h_{xx}&h_{xy}/\sqrt A&\mathsf L_x\\
 h_{xy}/\sqrt A&h_{yy}/A&\mathsf L_y/\sqrt A\\
 \mathsf L_x^T&\mathsf L_y^T/\sqrt A&\mathsf S
 \end{pmatrix}.                                       \tag{WA8}
\]
Block notation in (WA8) has sizes $1,1,3$. Every coefficient of
this finite polynomial matrix in $u,v$ is bounded and pure source.
The checker first expands the original signed core and all six
original event packets to verify the complete hard block and both
normal rows. It then verifies (WA5)--(WA7) from the corresponding
four plus two contact equations. Thus (WA8) is not merely a graph
or Hessian certificate.

\paragraph{Strict residue and a legal separated Gram lift.}
For infinitesimal independent $u,v$, the residue of (WA8), in the
printed root order, is
\[
 \overline{\mathsf G}=
 \begin{pmatrix}
 \bar n+\bar m+1/2&0&0&-\bar m&0\\
 0&1&0&0&0\\
 0&0&\bar C+\bar D+1/2&0&0\\
 -\bar m&0&0&\bar m&0\\
 0&0&0&0&\overline{\alpha/B^2}
 \end{pmatrix}\succ0.                                \tag{WA9}
\]
Indeed the nontrivial $2\times2$ determinant is
$\bar m(\bar n+1/2)>0$, and (WA2) supplies the last positive
pivot. The $q$ couplings disappear because $A^{-1/2}$ is
infinitesimal while $Ah$ is bounded. Every other omitted entry
has a factor $u,v$, or $uv$.

For completeness, strict bounded residue here licenses a finite
tensor SOS without a mixed Cholesky or inverse. Let
$\mathsf G_0=\mathsf G(0,0)$, a positive definite pure-source
matrix whose inverse is bounded. Whiten by its pure-source square
root and write the resulting matrix as
\[
 I+\sum_{j=1}^N H_j t_j,
\]
where each $H_j$ is a bounded pure-source symmetric matrix and
each $t_j$ is a nonconstant pure-testing monomial in $u,v$, hence
infinitesimal. Split $I$ into $N$ equal copies and orthogonally
diagonalize each $H_j$ purely in the source field. A resulting
scalar multiplier is $1/N+\ell t_j$, with $\ell$ bounded.
Choose a positive real $R>N|\ell|$ and put
$a=N\ell/R$, $b=Rt_j$. Then $|a|<1$, $|b|<1$, and
\[
 \frac1N+\ell t_j
 =\frac1{2N}\{(1+a)(1+b)+(1-a)(1-b)\}.               \tag{WA10}
\]
Each product is the square of a pure-source square root times
a pure-testing square root. This supplies finitely many legal
tensor roots for (WA8). Only finitely many source eigenvalues
occur, so one real $R$ can be chosen for all of them. If $N=0$
the source Gram itself is already sufficient.

This construction works for either sign and for zero $u$ or $v$,
and permits all relative source/testing rates. Applying the real
identity separately to the real and imaginary state parts proves
the complex-state assertion. It uses pure-testing square roots of
$1\pm Rt_j$, and consequently makes no claim to a global
polynomial SOS in the testing variables. There is therefore no
conflict with the obstruction proved on the physical SO1 family.

\paragraph{Exact scope of the gain.}
The result closes the local both-infinitesimal testing problem on
(WA1), with all four original state coordinates retained and with
no centered-packet subtraction. It does not yet cover the allowed
source openings $z/B$ infinitesimal or unbounded, or $1-r$
infinitesimal: in the latter case the $R_2$ pivot in (WA9) can
lose residue. A remote-testing theorem must still be joined before
this local certificate is called an all-tests source admission.
No assertion of nonemptiness of the final residual intersection
is made. The accompanying checker verifies thirteen exact
identities; the source bounds and finite lift above are the
positivity proof.
